% =============================================================================
\section{Limitations}
\label{sec:limitations}
% =============================================================================

We report our limitations with deliberate candor. Top venues reward honest
analysis; we believe documenting infrastructure failures alongside algorithmic
results is itself a contribution to reproducible science.

% -----------------------------------------------------------------------
\subsection{Infrastructure Failures}
\label{sec:infra_failures}
% -----------------------------------------------------------------------

\paragraph{JWT Token Expiration (Tinker API).}
Of 14 Tinker experiments launched, \textbf{11 failed mid-run} because the
platform's short-lived JWT authentication token expired before training
completed. Only three experiments ran to completion: DeepSeek-V3.1,
Qwen3-8B, and Llama-3.1-8B (tool\_use configuration). This is a fundamental
limitation of serverless ML platforms that do not support long-running
stateful jobs out of the box. Our workaround---issuing tokens immediately
before job submission---reduces but does not eliminate the hazard for
multi-hour runs. We recommend that platform operators expose a token-refresh
endpoint or implement background credential rotation; until then, practitioners
should checkpoint frequently and implement automatic restart logic. All
reported Tinker results in this paper derive from the three completed runs;
the remaining 11 experimental cells are treated as missing data and excluded
from aggregated statistics.

\paragraph{Modal Timeout (60-Minute Hard Limit).}
Three of six Modal-hosted evaluation jobs timed out at the platform's
60-minute wall-clock limit: held-out evaluations on 32B and 27B models, and
the HumanEval pass@1 suite. Large-model inference on a single A100 is
substantially slower than expected for generation-heavy tasks---generating
even modest numbers of long solutions at 32B scale can exceed 60 minutes
easily. Affected cells are marked \texttt{[TIMEOUT]} in the tables.
Future work should shard generation across multiple GPUs (tensor parallelism),
use speculative decoding, or negotiate higher time limits with the provider.

\paragraph{KL Divergence Tracking Bug.}
The gradient computation for KL divergence monitoring failed because reference
model logits were computed under \texttt{torch.no\_grad()}, but the downstream
KL term expected gradients from the reference branch. The resulting
\texttt{RuntimeError} (``element 0 of tensors does not require grad'') silently
killed the tracking loop without halting training, so policy-drift monitoring
was absent for all completed Tinker runs. Consequently, Table~\ref{tab:kl_divergence}
contains no Tinker-platform entries. The fix---retaining gradients for the
reference model or using \texttt{.detach()} consistently on both branches
before computing the KL estimate---is straightforward but was identified only
in post-hoc code review. We release the corrected implementation in
\texttt{src/grpo\_trainer.py} and note that this bug class is likely to affect
other GRPO re-implementations.

\paragraph{Weights \& Biases Step-Level Data Loss.}
All training runs completed W\&B sync without error, yet querying the API
for step-level reward trajectories returns zero steps for reward metrics across
all experiments---only run-level \emph{summary} values (final reward, episode
count) were captured. We traced this to a logging pattern that flushed metrics
to the summary dict (\texttt{wandb.run.summary.update(\ldots)}) rather than
calling \texttt{wandb.log(\{..., "step": step\})} at each training step,
causing W\&B to record a single aggregate point instead of a time series.
As a result, learning curves presented in Appendix~\ref{app:curves} are
reconstructed from local CSV checkpoints rather than W\&B trajectories, and
the published W\&B runs should \emph{not} be used to infer convergence speed.
The corrected logging code is included in the repository.

% -----------------------------------------------------------------------
\subsection{Methodological Limitations}
\label{sec:method_limits}
% -----------------------------------------------------------------------

\paragraph{Closed-Source Training Implementation (Tinker).}
Tinker is a commercial, closed-source API. We cannot inspect the exact GRPO
loss formulation, reward normalization scheme, minibatch construction, or
hardware configuration used server-side. Our Tinker results therefore measure
the \emph{platform's} GRPO implementation, not a precisely specified algorithm.
Researchers wishing to attribute performance differences to specific
implementation choices should use the open-source backends (TRL, OpenRLHF,
veRL) where every hyperparameter is auditable.

\paragraph{Short Training Horizons (30 Steps).}
All Tinker experiments used a budget of 30 gradient steps---a deliberate choice
to contain API costs but one that may be insufficient to observe convergence
on harder tasks. Thirty steps represents roughly one or two passes over the
prompt pool at the batch sizes we used. Long-horizon effects such as reward
hacking, catastrophic forgetting, or late-stage policy collapse are unlikely
to manifest at this scale. We regard our Tinker results as \emph{early-training
snapshots} rather than converged solutions, and caution against drawing strong
conclusions about asymptotic performance.

\paragraph{Single-Seed Tinker Experiments.}
Cost constraints precluded multi-seed replication on Tinker: each configuration
was run once. Without variance estimates we cannot apply standard significance
tests to Tinker results. We report these numbers descriptively and recommend
that future work with larger budgets run at least three seeds, consistent with
best practices from \citet{henderson2018deep}.

\paragraph{Train-Set Reward as the Primary Metric.}
Reported rewards are computed on the same prompt distribution used for
training, not a held-out test split. Only two partial held-out evaluations
were completed before timeout (Section~\ref{sec:infra_failures}). There is
therefore a non-trivial risk that reported gains partly reflect prompt-level
memorization rather than genuine generalization. Our open-source baselines
(TRL, CleanRL) do include held-out evaluation, providing a partial check;
but the Tinker-API results lack this safeguard.

\paragraph{Tool-Use Task: 0\% Reward.}
The synthetic tool-use task yielded a reward of exactly 0\% for all completed
runs under the Tinker backend. We interpret this as a task-design problem
rather than a model capability failure: the reward function requires exact
JSON-schema compliance with nested function signatures, a level of precision
that is difficult to reach in 30 steps without a warm-started SFT initialization.
The task may also require a graduated curriculum (simple → complex tool calls)
that was absent from our design. We release the task definition and reward
function so that the community can iterate; in the interim, tool-use results
should be treated as a negative baseline rather than a performance signal.

% =============================================================================
\subsection{Broader Impact}
\label{sec:impact}
% =============================================================================

\paragraph{Alignment and Misuse.}
GRPO and related RLHF methods are dual-use technologies. On the positive side,
they are the primary mechanism by which frontier models are aligned to human
preferences, and our benchmark lowers the barrier to studying these methods
rigorously. On the negative side, the same reward-maximization machinery can
be directed at adversarial objectives---optimizing for reward-hack proxies,
generating persuasive misinformation, or circumventing safety classifiers. We
do not provide new attack capabilities, but improvements in RL post-training
efficiency reduce the compute budget required for such misuse. We encourage
the community to develop robust reward modeling and out-of-distribution
detection in parallel with efficiency advances.

\paragraph{Democratizing RL-for-LLM Research.}
A primary motivation for \ours{} is to make rigorous RL post-training
experiments accessible to researchers without large-scale compute. By
publishing standardized tasks, Docker containers, and cost-annotated run
logs, we aim to reduce the advantage that well-resourced labs hold in this
space. Our total experimental expenditure was modest: approximately \$120 in
Tinker API credits (of which roughly \$80 was consumed by the 11 failed runs),
and roughly \$95 in Modal compute for the six held-out evaluation jobs
(three of which timed out). Open-source backend experiments were conducted on
a single A100 node donated by PES University's LLM Research Group. We include
these cost breakdowns explicitly so that other groups can budget accordingly.

\paragraph{Training Data and Privacy.}
All training data is drawn from publicly released datasets: GSM8K
\citep{cobbe2021gsm8k} for mathematical reasoning, HumanEval
\citep{chen2021codex} for code generation, and a synthetic tool-use corpus
generated without human subjects. No private data, personally identifiable
information, or licensed proprietary content is used. The benchmark does not
involve human participants in any capacity, and no IRB review was required.

\paragraph{No Direct Dual-Use Concern.}
Our contribution is a \emph{benchmarking framework}, not a new model or
capability. We do not release any model trained on sensitive, harmful, or
restricted content. All evaluated checkpoints are derived from publicly
available base models under permissive licenses (Apache 2.0 or equivalent).
We do not foresee direct security or biosafety risks arising from this work.

\paragraph{Reproducibility Statement.}
Despite the infrastructure failures documented above, all successfully
completed experiments are reproducible via Docker containers, fixed random
seeds, and step-by-step commands in \texttt{REPRODUCE.md}. Corrected logging
and KL-tracking code is included in the release. We have archived all local
CSV training logs so that the step-level reward trajectories excluded from
W\&B are nonetheless available to other researchers.
